% Standalone chunk A.5: Rank-adaptive corollary (POST-FIX, Round 4).
% Covers: Cor rank_adaptive (R4 statement with full Freedman tau_k,
%   explicit overline_Theta_k, a-priori Delta_sigma_max hypothesis),
%   proof via Weyl + shift correction, Rem shift_correction.
% Source: conf.tex lines ~461-485, 1858-1935 (post-R4).
% Shared preamble for standalone chunks. Do not compile directly; \input{} it.
\documentclass{article}
\usepackage[utf8]{inputenc}
\usepackage[T1]{fontenc}
\usepackage{lmodern}
\usepackage[margin=1in]{geometry}
\usepackage{amsmath,amssymb,amsfonts,amsthm,bm,mathtools}
\usepackage{enumitem}
\usepackage{microtype}
\usepackage{xcolor}
\usepackage{natbib}
\usepackage{booktabs}
\usepackage{hyperref}

\newcommand{\R}{\mathbb{R}}
\newcommand{\N}{\mathbb{N}}
\newcommand{\E}{\mathbb{E}}
\newcommand{\Prob}{\mathbb{P}}
\newcommand{\op}{\mathrm{op}}
\DeclareMathOperator{\sign}{sign}
\newcommand{\cA}{\mathcal{A}}
\newcommand{\cF}{\mathcal{F}}
\newcommand{\cH}{\mathcal{H}}
\newcommand{\cT}{\mathcal{T}}
\newcommand{\cI}{\mathcal{I}}
\newcommand{\cW}{\mathcal{W}}
\newcommand{\cTprobe}{\mathcal{T}_{\mathrm{probe}}}
\newcommand{\cE}{\mathcal{E}}
\newcommand{\norm}[1]{\left\lVert #1 \right\rVert}
\newcommand{\tilO}{\widetilde{\mathcal{O}}}
\newcommand{\DynReg}{\mathrm{DynReg}}
\newcommand{\rank}{\mathrm{rank}}
\newcommand{\range}{\mathrm{range}}
\DeclareMathOperator{\TopEig}{TopEig}
\DeclareMathOperator*{\argmax}{arg\,max}
\DeclareMathOperator{\tr}{tr}

\newtheorem{assumption}{Assumption}
\newtheorem{theorem}{Theorem}
\newtheorem{corollary}{Corollary}
\newtheorem{proposition}{Proposition}
\newtheorem{lemma}{Lemma}
\newtheorem{remark}{Remark}


\title{Chunk A.5: Rank-adaptive guarantee\\(post-fix, Round 4 verification)}
\author{(excerpt for standalone review)}
\date{}

\begin{document}
\maketitle

\textbf{Setup context.}
\begin{itemize}[leftmargin=*,itemsep=1pt,topsep=1pt]
\item Matrix Freedman (Thm.\ matrix\_bernstein\_conf) gives, w.p.\ $\ge 1-2\delta$,
$\|\widehat M_k-\bar M_k^{\mathrm{probe}}-\widetilde B\|_\op\le\tau_k$
with the full Freedman radius
$\tau_k:=2R_X\sqrt{\log(2d/\delta)/m_k}+2R_X\log(2d/\delta)/(3m_k)+\|\Theta_k\|_\op$
(square-root + linear + truncation-bias offset).
\item Under Gaussian probes, $\widetilde B=-\delta_\sigma/(d+2)\cdot I_d$
is a \emph{scaled identity} (Lem.\ G\_unbiased\_conf), with sign
\emph{opposite} to $\delta_\sigma$.
\item $\bar M_k^{\mathrm{probe}}$ has $r$ nonzero eigenvalues $\ge\lambda_{\min}$,
$d-r$ zeros (Prop.\ segment\_factorization\_conf).
\end{itemize}

\textbf{What changed since Round~3 (R4 delta, this chunk).}
\begin{itemize}[leftmargin=*,itemsep=1pt,topsep=1pt]
\item \textbf{Statement/proof $\tau_k^\star$ alignment (R3 remaining).}
The corollary statement now defines $\tau_k^\star$ using the \emph{full
Freedman radius} (including the linear term and a computable upper bound
$\overline\Theta_k\ge\|\Theta_k\|_\op$), matching the proof. No hidden
sub-leading absorption step.
\item \textbf{$\overline\Theta_k$ made explicit (R3 remaining).} The
corollary instantiates $\overline\Theta_k = Cd\sqrt{\delta/T}$ via
Prop.\ theta\_bound, so the threshold is fully algorithm-computable.
\item \textbf{A-priori $\Delta_{\sigma,\max}$ clarified (R3 remaining).}
The hypothesis says ``a priori known'' explicitly; a remark notes
that if $\Delta_{\sigma,\max}$ is instead a CI radius at confidence
$\ge 1-\delta''$, that $\delta''$ must enter the union-bound budget.
\item \textbf{Probability reparameterization made explicit (R3 remaining).}
The ``$2\delta$ per segment, union-bounded into the overall $\delta$-budget''
phrase is now accompanied by an explicit reparameterization
$\delta\mapsto\delta/(2K)$ at the top of the regret-chain proof.
\end{itemize}

\textbf{What changed since Round~2 (R3 delta, retained).}
\begin{itemize}[leftmargin=*,itemsep=1pt,topsep=1pt]
\item Knowability fix: R2's ``known $|\widetilde B|$'' replaced with
hypothesis ``known upper bound $\Delta_{\sigma,\max}\ge|\delta_\sigma|$''.
\item Probability fixed: $1-\delta\to 1-2\delta$ (Freedman + truncation
event from chunk A.2).
\item Sign rationale in Rem.\ shift\_correction now correct
($\delta_\sigma>0$ pulls true \emph{down}, not zeros up).
\end{itemize}

\section*{Statement}

\begin{corollary}[Rank-adaptive guarantee]
\label{cor:rank_adaptive}
Assume either (a) $\sigma_\varepsilon^2$ is known \emph{a priori} and
$\epsilon_\times=0$ (exact probes, so $\widetilde B=0$ and
$\Delta_{\sigma,\max}:=0$), or (b) isotropic Gaussian probes together
with an \emph{a priori} known upper bound
$\Delta_{\sigma,\max}\ge|\hat\sigma^2-\sigma_\varepsilon^2|$ on the
variance-plugin error (so $\|\widetilde B\|_\op\le\Delta_{\sigma,\max}/(d+2)$;
Lem.\ G\_unbiased\_conf). Define the \emph{full Freedman radius}
\[
  \tau_k := 2R_X\sqrt{\log(2d/\delta)/m_k}
  + \tfrac{2R_X\log(2d/\delta)}{3m_k}
  + \overline\Theta_k,
  \qquad
  \tau_k^\star := \tau_k + \Delta_{\sigma,\max}/(d+2),
\]
where $\overline\Theta_k\ge\|\Theta_k\|_\op$ is a computable upper
bound; the explicit choice $\overline\Theta_k = C_\Theta\cdot d\sqrt{\delta/T}$
works with the absolute constant $C_\Theta$ from the proof of
Prop.\ theta\_bound (App.~C chunk; explicit and independent of $T,d,\delta$). If the true rank $r$ is
unknown, thresholding $\widehat M_k$ at $2\tau_k^\star$ (computable
by the algorithm) recovers the true rank with probability
$\ge 1-2\delta$ whenever $\lambda_{\min}\ge 4\tau_k^\star$, and the
main bound of Thm.\ spsc\_regret applies unchanged (per-segment
failure $\le 2\delta$ is absorbed into the overall $\delta$-budget by
reparameterizing $\delta\mapsto\delta/(2K)$ at the top of the
regret-chain proof). \emph{Remark on implementability.} If instead
$\Delta_{\sigma,\max}$ is obtained from a sample-variance confidence
interval valid w.p.\ $\ge 1-\delta''$, add $\delta''$ to the union-bound
budget.
\end{corollary}

\section*{Proof}

By Theorem matrix\_bernstein\_conf, with probability $\ge 1-2\delta$
every eigenvalue of $\widehat M_k$ is within $\tau_k$ of the
corresponding eigenvalue of $\bar M_k^{\mathrm{probe}}+\widetilde B$,
where $\tau_k$ is the full Freedman radius defined in
Cor.\ rank\_adaptive (replacing $\|\Theta_k\|_\op$ by the computable
$\overline\Theta_k$ to match algorithm-side computation); Weyl gives
per-eigenvalue perturbation. Under the regime $m_k\ge\log(2d/\delta)$
and $\overline\Theta_k=O(d\sqrt{\delta/T})$ (Prop.\ theta\_bound) the
sub-leading terms are dominated by the square-root term---we keep them
explicit in $\tau_k$ so the algorithm-side threshold $2\tau_k^\star$
is uniformly valid, with no hidden absorption step.

Under Gaussian probes, $\widetilde B = -\delta_\sigma/(d+2)\cdot I_d$ is
a scaled identity (Lem.\ G\_unbiased\_conf), so the spectrum of
$\bar M_k^{\mathrm{probe}}+\widetilde B$ equals the spectrum of
$\bar M_k^{\mathrm{probe}}$ shifted by $-\delta_\sigma/(d+2)$.

Let $b:=\Delta_{\sigma,\max}/(d+2)$, a known upper bound on
$\|\widetilde B\|_\op$. Then the $r$ ``true'' eigenvalues of $\widehat M_k$
lie in $[\lambda_{\min}-b-\tau_k,\,\infty)$ and the $d-r$ ``zero''
eigenvalues lie in $[-b-\tau_k,\,b+\tau_k]$. Write
$\tau_k^\star := \tau_k+b$.

Under the eigengap hypothesis $\lambda_{\min}\ge 4\tau_k^\star$,
thresholding at $2\tau_k^\star$ separates the two groups:
``true'' $\ge\lambda_{\min}-\tau_k-b\ge 3\tau_k^\star > 2\tau_k^\star$
(kept); ``zero'' $\le\tau_k+b = \tau_k^\star < 2\tau_k^\star$
(discarded). Thresholding returns exactly $r$ indices with
probability $\ge 1-2\delta$.

Condition on this event; the regret analysis of Thm.\ spsc\_regret
applies verbatim with the estimated rank equal to $r$, so the main
bound holds, with the per-segment $2\delta$ cost absorbed by
reparameterizing $\delta\mapsto\delta/(2K)$ in the overall proof
union bound.\qed

\begin{remark}[Why the shift correction matters]
\label{rem:shift_correction}
The naive threshold $2\tau_k$ (without the $b$ correction) fails
under variance misspec. Since $\widetilde B = -\delta_\sigma/(d+2)\cdot I$
(Lem.\ G\_unbiased\_conf) has sign \emph{opposite} to $\delta_\sigma$:
if $\delta_\sigma > 0$, the shift $-\delta_\sigma/(d+2)\cdot I$ is
negative, pulling ``true'' eigenvalues \emph{down}---potentially
below the naive threshold; if $\delta_\sigma < 0$, the shift is
positive, pushing ``zero'' eigenvalues \emph{up}---potentially above
$2\tau_k$. Either way the rank estimate is wrong. The corrected
threshold $2\tau_k^\star = 2\tau_k+2b$ absorbs $\pm b$ shifts in
either direction; it is computable whenever $\Delta_{\sigma,\max}\ge|\delta_\sigma|$
is known. If $\sigma_\varepsilon^2$ is itself known a priori, take
$\delta_\sigma=0$ and the correction vanishes.
\end{remark}

\end{document}
